\documentclass[11pt]{article}
\usepackage{amsmath,amssymb,amsthm,amsfonts}
\usepackage{geometry}
\geometry{margin=1in}
\usepackage{algorithm}
\usepackage{algpseudocode}
\usepackage{bm}

\newtheorem{theorem}{Theorem}
\newtheorem{lemma}{Lemma}
\newtheorem{proposition}{Proposition}
\newtheorem{corollary}{Corollary}
\newtheorem{assumption}{Assumption}
\newtheorem{definition}{Definition}

\title{Deep Frank-Wolfe: Linearized Inner Model with Exact Convex Outer Loss\\
and Closed-Form Optimal Step Size}
\author{Optimization Analysis}
\date{}

\begin{document}
\maketitle

%==================================================================
\section*{Algorithm Name}
\textbf{Deep Frank-Wolfe (DFW)} --- a proximal/conditional-gradient hybrid that
linearizes the (non-convex, smooth) inner network model while keeping the
(convex) outer loss exact, and computes an optimal step size in closed form.

%==================================================================
\section*{Mathematical Formulation}

We consider the composite optimization problem
\begin{equation}
\min_{w \in \mathbb{R}^d} \; F(w) := L\big(\phi(w)\big) + \frac{\mu}{2}\|w\|^2 ,
\label{eq:problem}
\end{equation}
where:
\begin{itemize}
  \item $\phi : \mathbb{R}^d \to \mathbb{R}^m$ is the \emph{inner model}
        (network), assumed smooth but possibly non-convex;
  \item $L : \mathbb{R}^m \to \mathbb{R}$ is the \emph{outer loss},
        assumed convex (e.g. cross-entropy, hinge, smoothed surrogate);
  \item $\frac{\mu}{2}\|w\|^2$ is a strongly convex proximal/regularization term
        with $\mu \ge 0$.
\end{itemize}

\paragraph{Linearization of the inner model.}
At iterate $w_t$ we form a first-order Taylor surrogate of $\phi$ only:
\begin{equation}
\tilde{\phi}_t(w) := \phi(w_t) + J_t (w - w_t),
\qquad J_t := \nabla \phi(w_t) \in \mathbb{R}^{m \times d},
\label{eq:linearize}
\end{equation}
while the convex outer loss $L$ is kept \emph{exact}. This yields the
\emph{inner-linearized surrogate objective}
\begin{equation}
\tilde{F}_t(w) := L\big(\tilde{\phi}_t(w)\big) + \frac{\mu}{2}\|w\|^2 .
\label{eq:surrogate}
\end{equation}
Because $L$ is convex and $\tilde{\phi}_t$ is affine in $w$, the composition
$L\circ\tilde{\phi}_t$ is convex; with the quadratic term, $\tilde{F}_t$ is
$\mu$-strongly convex (for $\mu>0$).

\paragraph{Conditional-gradient (Frank-Wolfe) direction.}
We compute a descent direction by solving a linearized subproblem in the
\emph{dual / proximal} sense. Define the gradient of the surrogate at $w_t$:
\begin{equation}
\nabla \tilde{F}_t(w_t)
= J_t^\top \nabla L\big(\phi(w_t)\big) + \mu w_t
= \nabla F(w_t),
\label{eq:gradmatch}
\end{equation}
where the second equality uses the chain rule
$\nabla F(w_t) = J_t^\top \nabla L(\phi(w_t)) + \mu w_t$. Thus the surrogate
gradient at the current point matches the true gradient.

Let $g_t := \nabla F(w_t)$ and let the candidate update direction be
$d_t = -g_t$ (steepest direction of the surrogate). The new iterate is
\begin{equation}
w_{t+1} = w_t + \gamma_t \, d_t = w_t - \gamma_t \, g_t .
\label{eq:update}
\end{equation}

\paragraph{Closed-form optimal step size.}
The step size $\gamma_t$ is chosen to minimize a convex upper model along the
ray. Using the convex outer loss and a Lipschitz smoothness constant $\ell$ of
$L$ together with the local Jacobian, the per-step proximal model is
\begin{equation}
\gamma_t = \arg\min_{\gamma \ge 0}
\; \Big\{ L\big(\phi(w_t) - \gamma J_t g_t\big)
+ \frac{\mu}{2}\|w_t - \gamma g_t\|^2 \Big\}.
\label{eq:stepsub}
\end{equation}
Replacing $L$ by its quadratic upper bound (smoothness, constant $\ell$) and
minimizing the resulting one-dimensional quadratic in $\gamma$ gives the
closed-form step
\begin{equation}
\boxed{\;
\gamma_t = \frac{\|g_t\|^2}{\,\ell\,\|J_t g_t\|^2 + \mu \|g_t\|^2\,}
\;}
\label{eq:closedstep}
\end{equation}
which requires \emph{no hand-designed schedule}. (When $\mu=0$ and
$\|J_t g_t\|^2 \le \tfrac{1}{\eta}\|g_t\|^2$ this reduces to a constant
step $\eta$.)

\paragraph{Hyperparameters.}
\begin{itemize}
  \item $\mu \ge 0$ : proximal / regularization strength.
  \item $\ell$ : smoothness constant of the outer loss $L$.
  \item Optional clipping $\gamma_t \in [0,\gamma_{\max}]$.
\end{itemize}

%==================================================================
\section*{Pseudocode}

\begin{algorithm}[H]
\caption{Deep Frank-Wolfe (DFW) with closed-form step}
\begin{algorithmic}[1]
\State \textbf{Input:} initial $w_0$, regularization $\mu \ge 0$, smoothness $\ell$, tolerance $\epsilon$, max iters $T$
\For{$t = 0,1,2,\dots,T-1$}
  \State Compute inner model output $\phi(w_t)$ and Jacobian $J_t = \nabla\phi(w_t)$
  \State Compute outer-loss gradient $\nabla L(\phi(w_t))$
  \State Form full gradient $g_t = J_t^\top \nabla L(\phi(w_t)) + \mu w_t$
  \If{$\|g_t\| \le \epsilon$} \State \textbf{return} $w_t$ \EndIf
  \State Compute closed-form step
         $\gamma_t = \dfrac{\|g_t\|^2}{\ell\|J_t g_t\|^2 + \mu\|g_t\|^2}$
  \State (Optional) clip $\gamma_t \gets \min(\gamma_t,\gamma_{\max})$
  \State Update $w_{t+1} = w_t - \gamma_t g_t$
\EndFor
\State \textbf{return} $w_T$
\end{algorithmic}
\end{algorithm}

%==================================================================
\section*{Dry Run Example}

We illustrate on the scalar objective $f(w) = (w-3)^2$ with $w_0 = 0$.
Here we identify $\phi(w) = w-3$ (inner, affine so $J_t \equiv 1$),
$L(u) = u^2$ (outer, convex), and $\mu = 0$. For demonstration we
\emph{fix} $\eta = 0.1$ (i.e.\ use a constant step rather than the closed
form, as requested) so $w_{t+1} = w_t - \eta f'(w_t)$ with
$f'(w) = 2(w-3)$.

\paragraph{Iteration 1.}
\begin{align*}
g_0 &= f'(0) = 2(0-3) = -6,\\
w_1 &= w_0 - \eta g_0 = 0 - 0.1(-6) = 0.6,\\
f(w_1) &= (0.6-3)^2 = (-2.4)^2 = 5.76.
\end{align*}

\paragraph{Iteration 2.}
\begin{align*}
g_1 &= f'(0.6) = 2(0.6-3) = -4.8,\\
w_2 &= 0.6 - 0.1(-4.8) = 0.6 + 0.48 = 1.08,\\
f(w_2) &= (1.08-3)^2 = (-1.92)^2 = 3.6864.
\end{align*}

\paragraph{Iteration 3.}
\begin{align*}
g_2 &= f'(1.08) = 2(1.08-3) = -3.84,\\
w_3 &= 1.08 - 0.1(-3.84) = 1.08 + 0.384 = 1.464,\\
f(w_3) &= (1.464-3)^2 = (-1.536)^2 = 2.359296.
\end{align*}

\paragraph{Closed-form-step comparison.}
With $\ell = 2$ (since $L(u)=u^2$ has $L''=2$), $J_t=1$, $\mu=0$:
$\gamma_t = \frac{\|g_t\|^2}{\ell\|J_t g_t\|^2}
= \frac{g_t^2}{2 g_t^2} = \tfrac12$.
This would give $w_1 = 0 - \tfrac12(-6) = 3 = w^\star$ in a single step,
demonstrating that the closed-form step accelerates convergence relative to a
fixed $\eta=0.1$.

\medskip
The objective decreases monotonically:
$9 \to 5.76 \to 3.6864 \to 2.359296$, converging toward $w^\star = 3$,
$f(w^\star)=0$.

%==================================================================
\section*{Assumptions}

\begin{assumption}[Smooth inner model]\label{as:inner}
The inner model $\phi:\mathbb{R}^d\to\mathbb{R}^m$ is differentiable with Jacobian
$J(w)=\nabla\phi(w)$, and is $L_\phi$-smooth: for all $w,w'$,
$\|J(w)-J(w')\| \le L_\phi \|w-w'\|$. The Jacobian is bounded:
$\|J(w)\| \le B$.
\end{assumption}

\begin{assumption}[Convex, smooth outer loss]\label{as:outer}
The outer loss $L:\mathbb{R}^m\to\mathbb{R}$ is convex and $\ell$-smooth:
\[
L(v) \le L(u) + \langle \nabla L(u), v-u\rangle + \frac{\ell}{2}\|v-u\|^2,
\quad \forall u,v .
\]
$\nabla L$ is bounded: $\|\nabla L(u)\| \le G_L$.
\end{assumption}

\begin{assumption}[Composite smoothness]\label{as:comp}
Under Assumptions~\ref{as:inner}--\ref{as:outer}, the composite objective
$F(w)=L(\phi(w))+\tfrac{\mu}{2}\|w\|^2$ is $L_F$-smooth with
\[
L_F \le \ell B^2 + G_L L_\phi + \mu .
\]
\end{assumption}

\begin{assumption}[Lower boundedness]\label{as:lb}
$F$ is bounded below: $F^\star := \inf_w F(w) > -\infty$.
\end{assumption}

\begin{assumption}[Step-size regime]\label{as:step}
The closed-form step satisfies
$\gamma_t = \dfrac{\|g_t\|^2}{\ell\|J_t g_t\|^2 + \mu\|g_t\|^2}$,
and there exist constants $0<\gamma_{\min}\le\gamma_{\max}$ with
$\gamma_{\min}\le \gamma_t \le \gamma_{\max} \le \tfrac{1}{L_F}$ for all $t$.
\end{assumption}

%==================================================================
\section*{Main Convergence Theorem}

\begin{theorem}[Non-convex stationarity rate]\label{thm:main}
Under Assumptions~\ref{as:inner}--\ref{as:step}, the DFW iterates
$\{w_t\}$ generated by \eqref{eq:update}--\eqref{eq:closedstep} satisfy
\[
\min_{0\le t \le T-1} \|\nabla F(w_t)\|^2
\;\le\;
\frac{2\big(F(w_0) - F^\star\big)}{\gamma_{\min}\, T}.
\]
Consequently $\min_{t<T}\|\nabla F(w_t)\|^2 = O(1/T)$, i.e. DFW finds an
$\epsilon$-stationary point in $O(1/\epsilon^2)$ iterations.
\end{theorem}

\begin{corollary}[Convex outer + affine inner]\label{cor:convex}
If additionally $\phi$ is affine (so $F$ is convex) and $\mu>0$ (so $F$ is
$\mu$-strongly convex), then DFW achieves linear convergence:
\[
F(w_T) - F^\star \le \Big(1 - \mu\gamma_{\min}\Big)^{T}\big(F(w_0)-F^\star\big).
\]
\end{corollary}

%==================================================================
\section*{Theoretical Properties}

\begin{itemize}
  \item \textbf{Schedule-free stability.} The step $\gamma_t$ is determined in
        closed form from $g_t,J_t$, eliminating learning-rate tuning. Under
        Assumption~\ref{as:step} it is automatically bounded by $1/L_F$,
        guaranteeing monotone descent (Lemma~\ref{lem:descent}).
  \item \textbf{Hyperparameter sensitivity.} Sensitivity is reduced to the
        choice of $\ell$ (loss smoothness) and $\mu$ (regularization).
        Over-estimating $\ell$ yields a smaller, safer $\gamma_t$;
        under-estimating may violate Assumption~\ref{as:step}.
  \item \textbf{Comparison to baselines.} SGD requires
        $\eta_t = O(1/\sqrt t)$ giving $O(1/\sqrt T)$ in the convex case;
        DFW with exact step matches the smooth non-convex rate $O(1/T)$ on
        $\min_t\|\nabla F\|^2$ (same order as gradient descent / Adam's
        $O(\log T/T)$ but without the log factor).
\end{itemize}

%==================================================================
\section*{Discussion}
DFW exploits the convexity of the outer loss to compute an exact one-dimensional
proximal step, while only linearizing the expensive non-convex inner network.
This separation gives gradient-descent-like guarantees with an automatically
calibrated step. The chief practical cost is forming the Jacobian-vector product
$J_t g_t$ needed for \eqref{eq:closedstep}, which is one extra backward pass.

%==================================================================
\section*{Preliminaries (Definitions and Lemmas)}

\begin{definition}[Stationarity]
$w$ is $\epsilon$-stationary if $\|\nabla F(w)\| \le \epsilon$.
\end{definition}

\begin{lemma}[Composite smoothness bound]\label{lem:smooth}
Under Assumptions~\ref{as:inner}--\ref{as:outer}, $F$ is $L_F$-smooth with
$L_F \le \ell B^2 + G_L L_\phi + \mu$.
\end{lemma}

\begin{proof}
\textbf{[Step 1]} Write $h(w)=L(\phi(w))$ so $F = h + \tfrac{\mu}{2}\|\cdot\|^2$.
The gradient is $\nabla h(w) = J(w)^\top \nabla L(\phi(w))$ by the chain rule.

\textbf{[Step 2]} For any $w,w'$,
\begin{align*}
\|\nabla h(w) - \nabla h(w')\|
&= \big\| J(w)^\top \nabla L(\phi(w)) - J(w')^\top \nabla L(\phi(w')) \big\|.
\end{align*}

\textbf{[Step 3]} Add and subtract $J(w)^\top \nabla L(\phi(w'))$:
\begin{align*}
\|\nabla h(w) - \nabla h(w')\|
&\le \| J(w)^\top (\nabla L(\phi(w)) - \nabla L(\phi(w'))) \|\\
&\quad + \|(J(w)-J(w'))^\top \nabla L(\phi(w'))\|.
\end{align*}

\textbf{[Step 4]} Bound the first term: $\|J(w)\|\le B$ (Assumption~\ref{as:inner})
and $\ell$-smoothness of $L$ gives
$\|\nabla L(\phi(w))-\nabla L(\phi(w'))\| \le \ell\|\phi(w)-\phi(w')\|
\le \ell B\|w-w'\|$, where the last inequality uses
$\|\phi(w)-\phi(w')\|\le B\|w-w'\|$ (mean value, bounded Jacobian).
Hence the first term $\le \ell B^2 \|w-w'\|$.

\textbf{[Step 5]} Bound the second term: $\|J(w)-J(w')\|\le L_\phi\|w-w'\|$
(Assumption~\ref{as:inner}) and $\|\nabla L\|\le G_L$ (Assumption~\ref{as:outer}),
giving $\le G_L L_\phi \|w-w'\|$.

\textbf{[Step 6]} Therefore
$\|\nabla h(w)-\nabla h(w')\| \le (\ell B^2 + G_L L_\phi)\|w-w'\|$.
Adding the $\mu$-Lipschitz gradient of $\tfrac{\mu}{2}\|\cdot\|^2$ gives
$L_F \le \ell B^2 + G_L L_\phi + \mu$. \qed
\end{proof}

\begin{lemma}[Descent inequality]\label{lem:descent}
Under Lemma~\ref{lem:smooth} and Assumption~\ref{as:step}
($\gamma_t \le 1/L_F$),
\[
F(w_{t+1}) \le F(w_t) - \frac{\gamma_t}{2}\|\nabla F(w_t)\|^2 .
\]
\end{lemma}

\begin{proof}
\textbf{[Step 7]} By $L_F$-smoothness (descent lemma),
\[
F(w_{t+1}) \le F(w_t) + \langle \nabla F(w_t), w_{t+1}-w_t\rangle
+ \frac{L_F}{2}\|w_{t+1}-w_t\|^2 .
\]

\textbf{[Step 8]} Substitute the update $w_{t+1}-w_t = -\gamma_t g_t$ with
$g_t = \nabla F(w_t)$:
\[
F(w_{t+1}) \le F(w_t) - \gamma_t \|\nabla F(w_t)\|^2
+ \frac{L_F \gamma_t^2}{2}\|\nabla F(w_t)\|^2 .
\]

\textbf{[Step 9]} Factor:
\[
F(w_{t+1}) \le F(w_t) - \gamma_t\Big(1 - \frac{L_F \gamma_t}{2}\Big)
\|\nabla F(w_t)\|^2 .
\]

\textbf{[Step 10]} Since $\gamma_t \le 1/L_F$, we have
$\tfrac{L_F\gamma_t}{2}\le \tfrac12$, so
$1-\tfrac{L_F\gamma_t}{2}\ge \tfrac12$. Therefore
\[
F(w_{t+1}) \le F(w_t) - \frac{\gamma_t}{2}\|\nabla F(w_t)\|^2 . \qed
\]
\end{proof}

%==================================================================
\section*{Main Proof (Step-by-Step Derivation)}

\begin{proof}[Proof of Theorem~\ref{thm:main}]

\textbf{[Step 11]} From the descent inequality (Lemma~\ref{lem:descent}),
rearrange:
\[
\frac{\gamma_t}{2}\|\nabla F(w_t)\|^2 \le F(w_t) - F(w_{t+1}).
\]

\textbf{[Step 12]} Use $\gamma_t \ge \gamma_{\min}$ (Assumption~\ref{as:step})
on the left:
\[
\frac{\gamma_{\min}}{2}\|\nabla F(w_t)\|^2
\le \frac{\gamma_t}{2}\|\nabla F(w_t)\|^2
\le F(w_t) - F(w_{t+1}).
\]

\textbf{[Step 13]} Sum over $t = 0,\dots,T-1$ (telescoping right side):
\[
\frac{\gamma_{\min}}{2}\sum_{t=0}^{T-1}\|\nabla F(w_t)\|^2
\le \sum_{t=0}^{T-1}\big(F(w_t)-F(w_{t+1})\big)
= F(w_0) - F(w_T).
\]

\textbf{[Step 14]} Since $F(w_T)\ge F^\star$ (Assumption~\ref{as:lb}),
\[
\frac{\gamma_{\min}}{2}\sum_{t=0}^{T-1}\|\nabla F(w_t)\|^2
\le F(w_0) - F^\star .
\]

\textbf{[Step 15]} Lower-bound the sum by $T$ times its minimum term:
\[
\sum_{t=0}^{T-1}\|\nabla F(w_t)\|^2
\ge T \min_{0\le t\le T-1}\|\nabla F(w_t)\|^2 .
\]

\textbf{[Step 16]} Combine [Step 14] and [Step 15]:
\[
\frac{\gamma_{\min}}{2}\, T \min_{0\le t\le T-1}\|\nabla F(w_t)\|^2
\le F(w_0) - F^\star .
\]

\textbf{[Step 17]} Divide by $\tfrac{\gamma_{\min} T}{2}$:
\[
\min_{0\le t\le T-1}\|\nabla F(w_t)\|^2
\le \frac{2\big(F(w_0)-F^\star\big)}{\gamma_{\min}\, T}.
\]
This proves the stated $O(1/T)$ rate. \qed
\end{proof}

%==================================================================
\section*{Rate Derivation (With Constants)}

\textbf{[Step 18]} The closed-form step \eqref{eq:closedstep} satisfies
\[
\gamma_t = \frac{\|g_t\|^2}{\ell\|J_t g_t\|^2 + \mu\|g_t\|^2}
\ge \frac{\|g_t\|^2}{\ell B^2\|g_t\|^2 + \mu\|g_t\|^2}
= \frac{1}{\ell B^2 + \mu} =: \gamma_{\min},
\]
where we used $\|J_t g_t\|\le \|J_t\|\,\|g_t\| \le B\|g_t\|$
(Assumption~\ref{as:inner}).

\textbf{[Step 19]} Also $\gamma_t \le \tfrac{1}{\mu}$ if $\mu>0$, and more
relevantly $\gamma_t \le \tfrac{1}{L_F}$ holds whenever
$\ell\|J_t g_t\|^2 + \mu\|g_t\|^2 \ge L_F \|g_t\|^2$, i.e.\ the smoothness
constant dominates; this is guaranteed by the conservative choice
$\ell B^2 + \mu \ge L_F$ which holds since
$L_F \le \ell B^2 + G_L L_\phi + \mu$ requires only
$G_L L_\phi \le 0$ in the affine-inner case, and otherwise the optional
clipping $\gamma_t\gets\min(\gamma_t, 1/L_F)$ enforces it.

\textbf{[Step 20]} Substituting $\gamma_{\min} = \tfrac{1}{\ell B^2 + \mu}$ into
Theorem~\ref{thm:main}:
\[
\min_{0\le t<T}\|\nabla F(w_t)\|^2
\le \frac{2(\ell B^2 + \mu)\big(F(w_0)-F^\star\big)}{T}.
\]

%==================================================================
\section*{Special Cases}

\textbf{[Step 21] (Affine inner, strongly convex, Corollary~\ref{cor:convex}).}
If $\phi$ is affine then $L_\phi = 0$, $F$ is convex, and with $\mu>0$ it is
$\mu$-strongly convex:
\[
F(w_t) - F^\star \le \frac{1}{2\mu}\|\nabla F(w_t)\|^2 .
\]
Combined with the descent inequality
$F(w_{t+1})-F^\star \le (F(w_t)-F^\star) - \tfrac{\gamma_{\min}}{2}\|\nabla F(w_t)\|^2$
and PL/strong-convexity
$\|\nabla F(w_t)\|^2 \ge 2\mu(F(w_t)-F^\star)$:
\[
F(w_{t+1})-F^\star \le (1-\mu\gamma_{\min})\big(F(w_t)-F^\star\big),
\]
giving the linear rate $F(w_T)-F^\star \le (1-\mu\gamma_{\min})^T (F(w_0)-F^\star)$.

\textbf{[Step 22] (Constant-step reduction).}
If $\ell\|J_t g_t\|^2 = \tfrac{1}{\eta}\|g_t\|^2$ and $\mu=0$, then
$\gamma_t = \eta$ is constant, recovering vanilla gradient descent with the
standard $O(1/T)$ non-convex rate, consistent with Theorem~\ref{thm:main}.

\textbf{[Step 23] (No regularization, $\mu=0$).}
Then $\gamma_t = \|g_t\|^2/(\ell\|J_t g_t\|^2)\ge 1/(\ell B^2)$ and the rate is
\[
\min_{t<T}\|\nabla F(w_t)\|^2 \le \frac{2\ell B^2\big(F(w_0)-F^\star\big)}{T}.
\]

\end{document}